\chapter{周期律背后的拉锯战}

\begin{kaichang}
舀一小勺食盐，倒进一杯清水里，搅一搅。盐不见了——但你尝一口就知道它还在。前面几章讲过，食盐是钠离子 \chem{Na^+} 和氯离子 \chem{Cl^-} 整整齐齐堆成的晶体，下水之后，这些离子就散开，各自在水里游泳。

现在请你猜两件事。第一猜：一个钠原子（\chem{Na}）变成钠离子，丢掉了一个电子。它是变大了，还是变小了？第二猜：一个氯原子（\chem{Cl}）变成氯离子，多抢了一个电子。它是变胖了，还是变瘦了？

很多人凭直觉答：“丢东西的变小，得东西的变大。”这个答案对了一半，但错的地方更有趣——钠离子确实缩了水，可缩得远比你想的狠，半径几乎砍掉一半；氯离子也确实发了胖，胖得同样夸张。更奇怪的是，这一缩一胖，背后用的是同一套“力学”。

这套力学没有正式的名字，我们不妨叫它\textbf{拉锯战}：原子核拼命把电子往里拽，电子之间拼命把彼此往外推。第 10 章给了你周期表这张地图的轮廓，这一章的任务，是把地图上那些“往左变这样、往右变那样”的箭头趋势，一条一条从这场拉锯战里拆出来。拆完你会发现：周期律根本不用背，它只是两股力量记账的结果。
\end{kaichang}

\section{先给原子“量腰围”}

要谈拉锯战，先得知道战场有多大——原子到底多大？

这件事比量你的腰围麻烦得多。你可以拿软尺绕腰一圈，因为皮肤有一条明确的边界。原子没有。第 9 章讲过，原子核外的电子不是绕着固定轨道转的小行星，而是一团按概率分布的“电子云”，越往外越稀薄，却永远稀薄不到零。原子没有外壳，没有一条线可以说“到此为止”。

那书上那些原子半径的数字是怎么来的？办法很务实：不量单个原子，量\textbf{原子之间的距离}。比如在金属钠晶体里，用 X 射线衍射（一种靠波在规则排列上的干涉图案反推间距的技术）测出相邻两个钠原子核的距离，除以二，就当作钠原子的“半径”。在氯气分子 \chem{Cl_2} 里，测出两个氯原子核的距离除以二，又得到一个“半径”。两种量法的前提不同，给出的数字也不完全一样——所以严谨地说，原子半径不是一个绝对数字，而是“在某种相处方式下，原子愿意和别人保持多远”的度量。

先把“多远”翻译成直觉：原子的半径大约是 100 皮米（pm），也就是 $10^{-10}$ 米，一亿个原子肩并肩排起来，才大约 1 厘米长。你厨房里那卷铝箔，厚度大约 0.016 毫米，听起来薄得可怜——可除以铝原子约 $2.9\times 10^{-10}$ 米的直径，里面竟然叠着\textbf{五万多个原子}。

好，现在看测量结果。下面是第 2、3 周期一些元素的原子半径（不同量法数值略有出入，这里取一套常用值，单位 pm）：

\begin{table}[htbp]
\centering
\begin{tabular}{lccccccc}
\toprule
第 2 周期 & Li & Be & B & C & N & O & F \\
\midrule
半径/pm & 128 & 96 & 84 & 76 & 71 & 66 & 57 \\
\bottomrule
\end{tabular}
\end{table}

\begin{table}[htbp]
\centering
\begin{tabular}{lccccccc}
\toprule
第 3 周期 & Na & Mg & Al & Si & P & S & Cl \\
\midrule
半径/pm & 166 & 141 & 121 & 111 & 107 & 105 & 102 \\
\bottomrule
\end{tabular}
\end{table}

规律一眼就能读出来：\textbf{同一行从左往右，原子越缩越小}。再把锂（128）、钠（166）、钾（约 200 pm 出头）竖着比：\textbf{同一列从上往下，原子越来越大}。

奇怪吗？从左往右，电子明明越来越多，原子不胀反而缩；核的正电荷越来越多这件事，似乎比电子变多“更管用”。这就是拉锯战的第一条线索。我们该进场看看双方选手了。

\section{两股力量，一个战场}

把镜头推进到一个原子内部，清点一下场上的力量：

\begin{itemize}[itemsep=2pt]
\item \textbf{选手一：原子核的吸引力。}原子核带正电，电子带负电，异性相吸。这个吸引力有个要命的脾气：\textbf{对距离极其敏感}。距离拉大到两倍，吸引力就跌到四分之一（平方反比）；距离近一半，吸引力翻四倍。住在低楼层的电子被拽得死死的，住在高楼层的电子则明显被“手短”的原子核够得有点勉强。
\item \textbf{选手二：电子之间的排斥力。}电子都带负电，同性相斥。每个电子除了被核拽着，还被周围所有电子往外推。原子里的电子越多，这股“内讧”就越热闹。
\end{itemize}

一个原子的身材和脾气，就取决于这两位选手的记分牌。而记分牌上最关键的一个概念，叫\textbf{屏蔽效应}。

想象最外层的一个电子。它“抬头”看原子核，中间隔着内层一群电子。这群内层电子带负电，恰好挡在它和正电核之间，像一张网，把核的吸引力\textbf{挡掉了一部分}。于是外层电子实际感受到的，不是完整的核电荷，而是打了一个折扣的“净拉力”。这个折扣后的数字，叫做\textbf{有效核电荷}，常记作 $Z^{*}$。

你可以用一个简单的类比先感受一下这两股力量：

\begin{shiyan}
\textbf{静电里的吸引与屏蔽}

\textbf{材料}：一把塑料直尺（或气球）、干燥的头发或毛衣、一小堆碎纸屑（撕成米粒大小）、一张厨房铝箔。

\textbf{步骤}：
\begin{enumerate}[itemsep=1pt]
\item 把碎纸屑撒在干燥的桌面上。
\item 用尺子在头发或毛衣上快速摩擦十几下，然后慢慢靠近纸屑（不要碰到）。记下尺子离多远时纸屑开始“跳起来”被吸走。
\item 重复几次，故意改变距离，感受“近才吸得动”。
\item 让同伴把铝箔竖直立在尺子和纸屑之间（别让尺子碰到铝箔），再把摩擦过的尺子凑近——这次纸屑还动吗？移开铝箔再试一次。
\end{enumerate}

\textbf{观察点}：尺子越近，纸屑越容易被吸起，稍远一点就纹丝不动——吸引力随距离衰减得非常快。而一张薄薄的铝箔往中间一挡，纸屑立刻“装死”，吸引力像被关掉了。

\textbf{安全提示}：只允许用摩擦头发、毛衣的方式起电；\textbf{严禁}在插座、电线附近玩静电，更不许把任何金属物品插进插座。铝箔边缘锋利，小心割手。

\textbf{这个实验在模拟什么}：尺子上的静电像原子核的正电荷，距离越远越拉不动——这就是“选手一”的脾气；铝箔把电场挡了下来，这就是“屏蔽”的宏观版。当然，要说明白：原子里的屏蔽是电子云之间量子层面的排斥与排布，跟铝箔挡住电场只是\emph{类似}，不是同一机制——类比只负责帮你建立直觉。
\end{shiyan}

\section{一张心算记分牌：有效核电荷}

现在把“屏蔽”变成一个可以心算的粗糙模型，规则只有两条：

\begin{qiaoliang}
\textbf{有效核电荷的心算规则（粗糙但够用）}

把原子的电子分成“内层”和“最外层”。假设：每个内层电子大约抵消 1 份核电荷；同一层的电子彼此挤在一起，几乎挡不住同伴看核的视线。

于是：$Z^{*} \approx$ 核电荷数 $-$ 内层电子数，也就是约等于“最外层有几个电子”。

举几个例子：
\begin{itemize}[itemsep=1pt]
\item 锂（\chem{Li}，核电荷 +3，电子排布 2、1）：内层 2 个电子挡掉 2 份，最外层那 1 个电子感受到的净拉力约为 $3-2=1$。
\item 钠（\chem{Na}，+11，排布 2、8、1）：$11-10=1$。钾（\chem{K}，+19，排布 2、8、8、1）：$19-18=1$。——\textbf{同一族的有效核电荷几乎相同}。
\item 氟（\chem{F}，+9，排布 2、7）：最外层每个电子感受到的净拉力约为 $9-2=7$。氯（\chem{Cl}，+17，排布 2、8、7）：$17-10=7$。
\item 沿第 3 周期走：钠约 1，镁约 2，铝约 3……一直到氯约 7。\textbf{同一周期从左到右，净拉力一格一格地加码}。
\end{itemize}
\end{qiaoliang}

这张记分牌一立起来，周期表上两大半径趋势立刻有了因果解释：

\textbf{为什么同一周期从左往右原子缩小？}因为电子是加在\emph{同一层}上的，楼没盖高，可核的电荷在一格一格地加，而同层电子互相又挡不住——每个电子感受到的净拉力 $Z^{*}$ 从 1 涨到 7。拉得越狠，整团电子云被勒得越紧。从锂到氟，不是原子“变瘪了”，是拔河获胜的一方换成了原子核。

\textbf{为什么同一族从上往下原子变大？}因为 $Z^{*}$ 明明差不多（都是 1 或都是 7），电子却住进了更高一层楼。别忘了吸引力对距离极其敏感：楼层一高，距离一远，核就够不太着了。这一局，赢的是距离。

拉锯战的比分，就是这样写进周期表的。记住这个推理过程，比记住任何箭头方向都重要——\textbf{趋势不是背的，是算出来的}。

\section{抢走一个电子，要花多少钱}

半径描述的是战场的大小，接下来这个问题更直接：\textbf{把一个电子从原子里硬拽出来，要付多少“赎金”？}

这笔赎金有名字，叫\textbf{电离能}：把一个气态原子的最外层电子拿走、让它变成正离子，所需要输入的能量。用符号写就是：

\[\chem{Na(g)} \rightarrow \chem{Na^{+}(g)} + e^{-} \qquad \Delta E = +496\ \mathrm{kJ/mol}\]

注意是“$+$”496：能量是你\textbf{付出去}的，不是原子赏给你的。这呼应了全书反复强调的一条账规：拆散一对互相吸引的东西永远要耗能，让吸引双方结合才会放能——以后讲化学键时你会反复用到这条账规（第 17 章）。

下面是第 2 周期元素的第一电离能（单位 kJ/mol）：

\begin{table}[htbp]
\centering
\begin{tabular}{lcccccccc}
\toprule
元素 & Li & Be & B & C & N & O & F & Ne \\
\midrule
第一电离能 & 520 & 900 & 801 & 1086 & 1402 & 1314 & 1681 & 2081 \\
\bottomrule
\end{tabular}
\end{table}

大方向完全符合拉锯战的预言：从左到右，$Z^{*}$ 递增，电子被勒得越来越紧，赎金一路水涨船高，到最外层住满的氖达到顶峰。同族竖着看也吻合：钠 496，钾只要 419——钾的外层电子住四楼，离核远，自然好抢。这就是第 10 章那句“楼层越高越容易丢电子”的定量版本，也是第 1 族元素越往下脾气越爆的原因（它们遇水会怎样，第 12 章会细讲）。

但你有没有注意到表格里两个“打滑”的地方？铍（900）比硼（801）\emph{高}，氮（1402）比氧（1314）\emph{高}——明明往右走 $Z^{*}$ 在涨，赎金却不升反降。这不是测量错误，而是拉锯战里藏着的一条细则：\textbf{电子住的“户型”和“合住情况”也有发言权}。细节有点烧脑，放在本章末尾的挑战层，这里只记住结论：周期趋势是真实的大方向，但它是\textbf{有例外的趋势}，而例外本身也能用更深一层的能量记账解释——这恰恰是科学最迷人的样子。

再感受一下 496 kJ/mol 这个数量级。1 克金属钠大约是 $1/23$ 摩尔，把其中每个钠原子的外层电子全部抢走，需要的能量是：

\[496\ \mathrm{kJ/mol} \times \frac{1}{23}\ \mathrm{mol} \approx 21.6\ \mathrm{kJ}\]

21.6 kJ 是什么概念？水的比热大约是每千克每摄氏度 4.2 kJ，这点能量够把 1 千克水烧开约 $5\,^{\circ}\mathrm{C}$，或者让一只 60 瓦的灯泡亮 6 分钟。小小一克金属里，竟锁着这么一笔“电子赎金”——原子核的拉力，从来都不是闹着玩的。

有抢就有接。电子被拽出来后总得有去处，有些原子“接手”电子时不但不收费，还倒找能量：一个气态氯原子拿到一个电子变成氯离子，会\textbf{放出} 349 kJ/mol。这个数叫\textbf{电子亲和能}。放出能量说明什么？说明氯离子加自由电子的组合，总能量比氯原子加电子更低——还是那个原则：变化朝向能量更低的方向走，不是什么原子“想要”电子。

现在把两笔账合在一起看：钠丢电子要付 496，氯得电子只找回 349，净支出还有 147 kJ/mol。既然总账是亏的，食盐凭什么能稳定存在？先把这个疑问存在心里，章末的练习题会请你算这笔账——答案要等新的一股力量登场，那是第三篇的故事。

\section{电负性：拔河的比分牌}

电离能和电子亲和能说的都是“彻底抢走”或“彻底收下”。可更多时候，两个原子谁也不肯完全放手，而是\textbf{共享}一对电子，搭成化学键（第三篇的主线）。共享就公平吗？才不是。共享电子更像拔河绳中间系的那条红绸带——它几乎总是偏向力气大的一边。

物理学家给这个“拽共享电子的力气”起了名字：\textbf{电负性}。它是一张人为规定的比分牌：氟被定为全场最高，约 4.0；氧 3.4，氯 3.2，氮 3.0，碳 2.6，氢 2.2，钠只有 0.9，垫底的铯约 0.8。比分牌的走势你一定猜得到：往右上走，电负性一路涨到氟；往左下走，一路跌到铯。原因还是那两条：$Z^{*}$ 越大、楼层越低，核的手就越有力。

可这张比分牌是怎么定出来的？又不能拿弹簧秤去钩电子。这背后有一个很漂亮的证据故事。

\begin{zhengju}
1932 年，美国化学家鲍林（Linus Pauling）盯着一个看似琐碎的账目问题：为什么氢氯键 \chem{H-Cl} 的键能，比氢氢键 \chem{H-H} 和氯氯键 \chem{Cl-Cl} 键能的\textbf{平均值}大出一截？

他的推理是这样的。如果氢和氯对共享电子的拉力半斤八两，电子云均匀铺在中间，那么“混合键”的强度没有理由超过两种“纯键”的平均——拔河绳不偏不倚，谁也没占到便宜。可实验数据一次次显示：不同原子之间成键，几乎总是比平均值\textbf{更牢}。多出来的这份牢固从哪来？鲍林判断：电子被拉偏了。偏向一方之后，分子一端略负、一端略正，正负相吸，凭空多了一份静电的“奖金”，键才格外牢。拉得越偏，奖金越大。

于是他做了一件化学家很少干的事：把“超出平均值的那部分键能”换算成一张 0 到 4 的排行榜，宣布这就是各种原子拽电子的力气——电负性标度。整张表只从量热数据里算出来，没有直接“看”过任何一个电子。

这个定义当时就有争议：它毕竟是\textbf{反推}出来的间接量，不是天平上称出来的。所以后来不断有人提出替代算法——马利肯（Mulliken）说，拽电子的力气应该等于“抢电子的本事”和“留电子的本事”的平均，也就是电离能和电子亲和能取平均；奥尔雷德和罗周（Allred 与 Rochow）则直接从有效核电荷除以半径平方去估。几张排行榜细节互不相同，可名次几乎一模一样：氟永远是第一，铯永远垫底。几种互相独立的算法指向同一个顺序——这反而说明电负性抓到了某种真实的东西。科学里很多有用的概念都是这样：不能直读，但怎么算都跑不掉。
\end{zhengju}

电负性的用处大到超出你想象。两个原子电负性差得多，电子就严重偏向一边，形成食盐那种正负离子；差得少，电子勉强共享，形成水、糖、油脂里的共价键。水为什么能溶解盐却溶不了油？谜底就藏在电负性差造成的“分子正负极”里——这是第三篇要展开的故事。医院检验科和反兴奋剂实验室常用的质谱仪，第一步也是用电离能把样品分子打成带电离子，再按质量分选，连十亿分之一量级的违禁物质都逃不掉；街边的霓虹灯，则是靠高压把稀有气体原子的电子打掉、让气体导电发光——选稀有气体，正因为它们平时电负性意义上的“拔河”一概不参加，灯管才不会被腐蚀。

\section{趋势就藏在你的生活里}

这些“拉锯比分”不是实验室里的摆设，你身边的很多常识，其实都是它的推论。

生日派对上的氦气球，为什么敢放心地飘在孩子头顶？因为氦的外层电子赎金全场最贵之一，谁也不够钱抢，它也就什么反应都不参加。换成氢气球试试——氢气遇火会猛烈燃烧，所以大型活动的气球严格规定只许充氦气。化学实验室里，金属铯要密封在充了保护气的玻璃管里，隔着管子看它都金光闪闪地发软；它的电离能全场垫底，外层电子几乎是白送的，暴露在空气里瞬间就“成交”——它和水见面的场面有多火爆，第 12 章细讲。你长辈手指上的金戒指，戴几十年也不锈不暗，原因之一是金的电离能很高、电负性也不低，电子抓得牢，氧气很难跟它成交；而铁钉几个月就锈迹斑斑——金属之间的这份差别，连同铁锈的攻防，第 15 章会算总账。再看你体检报告里的血钾、血钠指标：它们测的其实都是 \chem{K^+} 和 \chem{Na^+}，因为这两种原子在体内早就把外层电子“脱手”了——赎金太便宜，根本留不住。

一句话：哪家的电子好抢、哪家的电子难动，整张周期表的表情，都被这场拉锯战提前写好了。

\section{这张地图的边界在哪里}

任何一个好用的模型都有失效的角落，拉锯战模型也不例外。用之前，先看清它的三包条款。

\begin{itemize}[itemsep=2pt]
\item \textbf{趋势有例外。}电离能表格里铍和氮的“打滑”已经见过：当电子恰好把某种户型住满或住到一半时，排布会额外安稳一点，赎金就额外贵一点。大趋势是坡度，例外是台阶——两者都是真的。
\item \textbf{中间街区不太听话。}在周期表中央的过渡金属区，电子填的是靠内的 d 层，屏蔽情况复杂，邻居之间的半径、电离能差别很小，趋势曲线几乎躺平。这正好解释了为什么那一区元素性格相近、又都能变出好几种价态——第 15 章专门逛这个街区。
\item \textbf{这些数字都不是“称”出来的。}半径依赖测量场合（金属里、分子里、晶体里数值各不同），电负性依赖定义方案。换人一套算法，数字就挪一点——但\textbf{趋势和排序稳如泰山}。用这些数字比大小、判方向，靠谱；拿小数点后一位较劲，就是缘木求鱼。
\item \textbf{离子的胖瘦还看邻居。}一个离子在晶体里被几个异号离子围着，胖瘦也会略有不同。所以离子半径同样是“在某种相处方式下”的数字。
\end{itemize}

一句话总结三包条款：拉锯战模型是\textbf{推理的脚手架}，不是背诵的清单。每当你想用一条趋势，先问一句：这一局里，$Z^{*}$、楼层、户型，各站在哪一边？

\begin{wuqu}
“电子那么轻，得失电子只是换个电荷符号，原子的大小基本不变。”

错得很有迷惑性，因为前半句是对的：一个电子的质量只有质子的约 1/1836，钠原子丢一个电子，体重几乎一两没掉。可决定原子胖瘦的不是电子的\emph{质量}，而是\emph{拉力和排斥的平衡}。钠原子半径约 186 pm，钠离子只剩约 102 pm——半径打了五五折，体积按立方算下来缩到原来的约 $(102/186)^3 \approx 1/6$！反过来，氯原子约 102 pm，氯离子却胖到约 181 pm。得失一个“几乎没有质量”的电子，竟能让原子的体积翻几倍或者缩到零头。下次再有人拿“电子很轻”论证“离子和原子差不多大”，就把这组数字拍给他。
\end{wuqu}

\section{回到那杯盐水}

现在可以兑现开场的两个谜题了。

\textbf{钠为什么缩水缩得那么狠？}钠原子的排布是 2、8、1，丢掉的是第三层那唯一的电子——整层楼被拆了，剩下的 10 个电子全挤在一、二层。而且电离之后核电荷仍是 +11，电子只剩 10 个，每个电子感受到的净拉力反而变大，电子云被整体勒紧。拆楼加收紧，双管齐下，半径从约 186 pm 直降到约 102 pm。

\textbf{氯为什么发胖？}氯原子排布 2、8、7，核 +17。抢到第 18 个电子之后，楼层没变，可核要拽的电子多了一个，电子之间的排斥也添了一份，每个电子分到的有效拉力被摊薄，整团电子云就松垮下来，从约 102 pm 涨到约 181 pm。

所以你面前这杯咸水里游来游去的，其实是一对“反差组合”：小个子的钠离子和大个子的氯离子。顺带说一句，它们还不是裸泳——每个离子身上都裹着一层水分子凑成的“外套”。水分子为什么心甘情愿排队给离子当外套？因为水分子自己有正极和负极——这又是电负性写下的伏笔，等到第三篇讲分子极性时再揭晓。

\section{拉锯战还远没有打完}

这一章我们把周期律拆成了两股力的账：核的吸引对距离的敏感，电子的排斥与屏蔽。半径、电离能、电子亲和能、电负性，四张表其实是同一张记分牌的四种读法。

接下来的两章，我们要去看这场拉锯战的两个极端赛区。第 12 章去最左边：第 1、2 族的元素把“丢电子”做到了极致，赎金低到一遇水就当场成交，那场面相当火爆。第 13 章去最右边：卤素把“抢电子”做到了极致，而稀有气体干脆退赛——不过“退赛”也不是绝对的，氙就被人硬是拉上过赛场。两股力量在两端的表演，比中间地带精彩得多。

\begin{tiaozhan}
\textbf{电离能为什么会“打滑”？}（跳过不影响主线）

细看第 2 周期：铍（900）比硼（801）高，氮（1402）比氧（1314）高。秘密在楼层内部的“户型”和“合住”。同一层里，s 户型比 p 户型更贴近核（量子力学叫钻穿效应），所以硼那个孤零零的 2p 电子，其实比铍那对贴核的 2s 电子更容易被拿走——赎金自然便宜了。氧的情况类似：2p 户型有三间“房”，按规则电子先一人一间（半满，三间各住一人），第 4 个电子被迫与其中一人合住。同处一房的两个电子排斥格外强，于是氧有一个电子“住得憋屈”，比氮的半满电子更容易被请走。所谓“全满、半满格外稳定”，不是神秘定律，而是\textbf{排斥最小化加户型对称}共同记出的账。精细的屏蔽算法（斯莱特规则）会给同层电子也记上零点几份的遮挡，结论与我们的心算模型一致，只是比分更准。
\end{tiaozhan}

\section*{练一练}

\begin{sikaoti}
\item 用心算规则估算镁（\chem{Mg}，核电荷 +12）和硫（\chem{S}，核电荷 +16）的最外层电子各自感受到的有效核电荷。预测：两者相遇，谁更可能交出电子、谁更可能收下电子？生成的离子各带什么电荷？
\item 画两幅“拉锯战”示意图：左边画氟原子（核 +9，9 个电子），右边画氟离子（核 +9，10 个电子）。用箭头标出核的拉力和电子间的排斥，并注明哪一个里“每个电子分到的拉力”更小。图旁注明：圆圈和箭头是人为的表示法，电子并不真的排成圆圈。
\item 把 \chem{Na}、\chem{Na^{+}}、\chem{Cl}、\chem{Cl^{-}} 四种粒子按半径从大到小排序，并用“楼层 + 屏蔽 + 排斥”三句话说明理由。
\item 画一张能量账图：横轴是反应进程，纵轴是能量。先画出“钠原子电离”这段上坡（$+496$ kJ/mol），再画出“氯原子得电子”这段下坡（$-349$ kJ/mol），标出净差 $+147$ kJ/mol。既然单独算总账是亏的，食盐晶体凭什么还能稳定存在？猜一猜还缺哪一笔账（提示：正负离子堆成晶体时会怎样），第 18 章对答案。
\item 第 2 周期里，铍的电离能比硼高、氮比氧高。用挑战层里“户型”和“合住”的语言，向同桌解释这两个例外为什么不违背周期趋势。
\item 把 2 克金属钠（摩尔质量约 23 g/mol）的每个原子都电离一次，大约需要多少能量？够把 1 千克水烧开几摄氏度？（水的比热约 4.2 kJ/(kg·$^{\circ}$C)。）
\end{sikaoti}
